\section{Introduction}
Arterial stiffness is a well-established early predictor for cardiovascular risk. In particular, increased aortic stiffness has been linked with numerous cardiovascular diseases such as stroke, hypertension, and coronary artery disease~\cite{sutton-tyrrellElevatedAorticPulse2005}. Aging is strong contributor to arterial stiffening, largely due to biochemical and structural alterations that thicken the vessel wall and reduce the aorta's ability to buffer pulsatile blood flows~\cite{vatnerVascularStiffnessAging2021,cavalcanteAorticStiffness2011}. Recent studies have indicated that increased aortic stiffness may also be linked to cranial pathologies such as vascular dementia and Alzheimer's disease~\cite{deroosMagneticResonanceImaging2017,batemanPulseWaveEncephalopathy2004}. Many of the risk factors associated with neurodegeneration such as smoking, hypertension, and the apolipoprotein E4 (APOE4) gene are also associated with cardiovascular disease~\cite{hughesReviewPotentialRole2015,nordestgaardSharedRiskFactors2022,angoffAorticStiffnessEpidemiology2021}. As such, there is high clinical interest in accurately assessing arterial stiffness. 

Direct measurement of arterial stiffness is challenging; however, reductions in arterial compliance manifest as increases in blood pulse wave velocity (PWV) as described by the Moens-Korteweg equation, providing a practical indirect metric~\cite{callaghanRelationshipPulsewaveVelocity1986}. The gold standard for PWV measurements is the use of pressure wires, but that is a highly invasive and uncomfortable approach which necessitates the need for non-invasive approaches. Most non-invasive methods utilize time-shift methods to determine regional PWV where flow or pressure waveforms are measured at multiple locations along the vessel to calculate vascular path lengths ($\Delta d$) and the time shifts ($\Delta t$) between them which gives $PWV=\frac{\Delta d}{\Delta t}$~\cite{wentlandReviewMRIbasedMeasurements2014}. 

There are multiple approaches to measuring time-shifts non-invasively such applanation tonometry and ultrasound, however, these methods utilize external distance measurements which do not take into account vessel tortuosity leading to incorrect PWVs~\cite{pereiraNovelMethodsPulse2015}. In recent years, MRI has emerged as a promising avenue for non-invasively measuring time-shifts through 2D phase contrast (2DPC) imaging which enables quantification of blood velocity and flow through a plane~\cite{wentlandReviewMRIbasedMeasurements2014}. These can be combined with standard anatomical images for accurate vessel path lengths leading to more robust PWV calculations. However, standard cine 2DPC cardiac imaging is acquired with Cartesian sequences which require breath holds to avoid repiratory motion artifacts. However, many elderly, senile, or health compromised patients such as those with Alzheimer's are incapable of holding their breath or following breathing directions which can result in severely compromised image quality. Alternative imaging methods with more flexible breathing requirements would thus be desirable. One such method would be radially sampled MRI sequences. They offer robustness to motion artifacts~\cite{gloverProjectionReconstructionTechniques1992}, flexible tradeoffs for spatial or temporal resolution, favorable noise properties for regularized reconstructions~\cite{lustigSparseMRIApplication2007}, and, importantly, more flexibility in physiological gating. By utilizing respiratory bellows to track respiratory motion, radial spokes acquired only during a static respiratory state (e.g., expiration) can be retained and reconstructed, allowing for a scan to be performed under free-breathing conditions. In addition to breath-holds, temporal resolution is also a limitation of MRI-based PWV\@. Higher temporal resolutions are necessary for capturing faster transit times that may result from age- or disease-related aortic stiffening. 

The primary aim of this work is to assess the feasibility of a free-breathing, radially-undersampled 2DPC sequence by comparing PWV measures to a standard, Cartesian 2DPC breath-hold sequence. This extends upon Radial data was reconstructed with a standard reconstruction (matching Cartesian image parameters) and a locally low rank reconstruction with increased temporal resolution. Aortic PWV values were obtained in a cohort (N=83) of middle-aged, cognitively unimpaired adults to determine intra- and inter-observer repeatability, establish normative PWV values, and assess the relationship between aortic PWV and age. Lastly, due to a lack of consensus in the optimal time shift method to use for MRI-based PWV measurements, 4 different transit time algorithms were employed to assess their impact on calculated PWVs.

\section{Materials \& Methods}
\subsubsection{Study Population}
This single-center, HIPAA-compliant study was approved by the UW-Madison Health Sciences Institutional Review Board and written informed consent was obtained for all participants. Data for this investigation came from 83 participants enrolled in the Longitudinal Impact of Fitness and Exercise (LIFE) study, drawn from the larger Wisconsin Alzheimer's Disease Research Center (ADRC) Clinical Core and Wisconsin Registry for Alzheimer's Prevention (WRAP) longitudinal, observational studies. These cohorts consist of healthy middle-aged control subjects and other older individuals at-risk for Alzheimer's disease (AD), including higher prevalence of APOE4 genotype and parental history of AD\@. Participants were included in the LIFE study if at least 45 years of age and if there were no signs of cognitive impairment or impaired decision making at the time of enrollment, determined via neuropsychological examinations~\cite{jackNIAAAResearchFramework2018,langhoughkoscikValidityEvidenceResearch2021}. Participants were excluded from the LIFE study if taking antipsychotic medications at the time of appointment or if MRI was generally contraindicated (e.g., pregnancy, magnetically-activated implant, etc.). MRI data for this study were obtained on follow up visits between January 2023 and December 2024. 

\subsubsection{MR Imaging Protocol}
Data were acquired on a 3T Signa Premier system with a 30-channel AIR™ anterior array coil (MR750, GE Healthcare, Waukesha, WI). For physiological gating, photoplethysmography (PPG) data was recorded during the scan using a peripheral pulse oximeter attached to the index finger and respiratory motion was recorded using respiratory bellows attached to the chest. 

An ungated, free-breathing, balanced steady-state free-precession (bSSFP) product sequence was used to attain anatomical images to guide 2DPC plane prescription and for aortic centerline tracing and distance measurements. Axial, coronal, and sagittal views were obtained with full aortic coverage using the following scan parameters: TR/TE = 3.0/1.3 ms; flip = 45$^\circ$; acquired FOV = 480 mm (axial scan = 360 mm); acquired matrix size = 148 $\times$ 148; reconstructed matrix size = 512 $\times$ 512; reconstructed spatial resolution = 0.94 $\times$ \qty{0.94}{\mm^2} (axial scan = 0.70 $\times$ \qty{0.70}{\mm^2}); slice thickness = 8 mm (sagittal scan = 3.2 mm); scan time = 8 \- 13 s. Local gradient shims were used to improve field homogeneity over the aorta and reduce banding artifacts. 

For both Cartesian and radial 2DPC sequences, two axial planes were placed in the aorta. The first slice was prescribed in the aortic arch, transecting both the ascending and proximal descending aorta. The second slice was prescribed just above the kidneys, transecting the proximal abdominal aorta (Figure 1A). This allowed for a total of 3 flow measurements in the aorta. Cartesian 2DPC scans (product sequence) were acquired with PPG gating during breath-hold on expiration with the following parameters: TR/TE = 5.1/3.1 ms; flip angle = 25$^\circ$; VENC = 150 cm/s; velocity encoding scheme = 2-point through-plane; views per segment (VPS) = 4; acquired FOV = 36 cm (60\% phase FOV); acquired matrix size = 192 $\times$ 160; acquired temporal resolution = 40.7 ms (defined as 2*TR*VPS); reconstructed matrix size = 256 $\times$ 256; reconstructed spatial resolution = 1.41 $\times$ \qty{1.41}{\mm^2}; slice thickness = 6 mm; reconstructed cardiac frames = 40; scan time = 10 \- 20 s (depending on heart rate). The acquired and standard reconstructed number of cardiac time frames were selected based on the publication from Dorniak et al.~\cite{dorniakRequiredTemporalResolution2016} on the required temporal resolution for 2DPC MRI\@. 

\begin{figure}[ht]
    \vspace{0.25cm}
    \centering
    \includegraphics[width=\textwidth]{PWV_tool.png}

    \caption[PWV Analysis tool]{A. Two axial planes were acquired for each PC scan, one bisecting the aortic arch and the other in the abdomen, for a total of three aortic ROIs. B. A custom software platform was developed to segment the aorta ROIs, trace aortic centerlines, analyze flow curves, and calculate pulse wave velocity.}\label{fig:pwv_tool}
    \vspace{0.25cm}
\end{figure}

Radial 2DPC scans (research sequence) were acquired with retrospective PPG cardiac and retrospective respiratory gating with the following parameters: TR/TE = 7.5/4.3 ms; flip = 25$^\circ$; VENC = 150 cm/s; velocity encoding scheme = 2-point through-plane; number of projections = 10,000; golden-angle sampling; acquired FOV = 32 cm; slice thickness = 6 mm; scan time = 2:30. For the radial sequence, two different types of reconstructions were employed including (1) a standard reconstruction designed to match spatial/temporal parameters to the Cartesian acquisition and (2) a high-temporal resolution reconstruction using a locally low rank (LLR) constrained reconstruction~\cite{jimenezFeasibilityHighSpatiotemporal2018,haldarLowrankApproximationsDynamic2011} with both reconstructions using ESPIRiT~\cite{ueckerESPIRiTEigenvalueApproach2014} coil sensitivity estimation, Maxwell term phase offset corrections~\cite{bernsteinConcomitantGradientTerms1998}, and gradient calibration to correct trajectory errors~\cite{johnsonImproved3DPhase2008}. The standard reconstruction was performed using Parallel Imaging with Localized Sensitivities (PILS) reconstruction to reduce undersampling artifacts~\cite{griswoldPartiallyParallelImaging2000} to obtain the following image parameters: matrix size = 256 $\times$ 256; spatial resolution = 1.40 $\times$ \qty{1.40}{\mm^2}; reconstructed cardiac frames = 40. The LLR reconstruction was performed with nuclear norm regularization using a local block size of 4 $\times$ 4 $\times$ 4, empirically tuned regularization parameter of $\lambda$ = 0.003, and 55 iterations to obtain the following image parameters: matrix size = 320 $\times$ 320; spatial resolution = 1.00 $\times$ \qty{1.00}{\mm^2}; reconstructed cardiac frames = 80. Time-resolved and time-averaged velocity and magnitude datasets were obtained for both reconstructions. All data were reconstructed in expiratory phase as determined by the respiratory bellow signal with a respiratory acceptance window of 50\%. 

\subsubsection{Image Post-Processing}
Quality assurance was performed in MATLAB (version 2020b, Mathworks, Natick, MA). All cardiac images were assessed for general image quality, proper plane placement, and velocity aliasing by viewing the time-resolved velocity data. Respiratory and cardiac gating was assessed using a customized MATLAB tool designed to read raw gating data, displaying both respiratory waveforms and recorded R-R intervals across the duration of the scan. Datasets with poor image quality, plane placement, gating quality, or velocity aliasing were rejected. 

In order to process cardiac 2DPC MRI data and obtain aortic PWV measures, a graphical user interface (GUI) was developed in MATLAB (Figure 1B) and is available online~\cite{roberts2DPhaseContrast}. The first component of the GUI was designed for centerline tracing from anatomical images to calculate distances ($\Delta d$) between ROI measurements (Figure 2). Axial, coronal, and sagittal bSSFP images were loaded separately into the tool. For each set of images, fiducial points were manually placed over the aorta throughout all slices. Next, the aortic arch and abdominal aorta 2DPC images were loaded into the tool and a point was placed in the center of the ascending, descending, and abdominal aorta (where the flow measurements were taken). All points were spatially localized by converting image coordinates to a physical coordinate system (relative to scanner isocenter) using DICOM header information. Following this, sagittal and coronal projections of these points were created, and the user was asked to trace a centerline through these set of fiducial points. The centerline was then interpolated to 150 points using modified Akima interpolation. Distance calculations between flow measurement locations were then automatically computed. 

\begin{figure}[ht]
    \vspace{0.25cm}
    \centering
    \includegraphics[width=\textwidth]{centerline.png}

    \caption[Centerline creation]{Free-breathing axial, coronal, and sagittal bSSFP images were acquired to trace aortic centerlines for path length calculations. The developed GUI walks users through the process by prompting them to place fiducial points identifying the aorta in the axial and sagittal views before tracing centerlines through the selected points. The centerlines are localized spatially and aligned with the acquired axial slices.}\label{fig:centerline}
    \vspace{0.25cm}
\end{figure}

The second component of the GUI was designed to calculate time shifts from flow waveforms and to calculate PWV values. First, the centerline data (obtained from the previous step) was loaded into the tool. Next, aortic arch and abdominal aorta 2DPC images were loaded and the time-averaged magnitude image was displayed. Three circular ROIs were placed within the ascending aorta, descending aorta, and abdominal aorta. The ROI was then propagated over all cardiac time frames and velocity waveforms were constructed by taking the average velocity within the ROI time frame. Raw velocity waveforms were initially inspected to assess waveform noise before interpolation and smoothing. Due to the use of PPG gating, the velocity waveforms were significantly shifted such that the systolic peaks were often located near the end of the constructed waveform. This physiological delay is caused by the pulse pressure traveling from the heart to the left index finger, where the pulse oximeter was located. Raw waveforms were shifted by a user-defined number of frames such that systole occurred in the first part of the cardiac cycle. From this, velocity waveforms were up-sampled to 1000 points and were subsequently smoothed using a Gaussian-weighted moving average kernel (window width = 80 interpolated sample points) to decrease velocity noise. Several commonly-used transit time algorithms (Figure 3) were incorporated into the tool to automatically determine temporal shift in velocity waveforms, including: time-to-point (TTP)~\cite{groeninkBiophysicalPropertiesNormalsized2001,doguiMeasurementAorticArch2011}, time-to-foot (TTF) with baseline offset correction~\cite{lundstromPulseWaveVelocity2022}, time-to-upstroke (TTU)~\cite{doguiMeasurementAorticArch2011}, and cross correlation (XCOR)~\cite{fieldenNewMethodDetermination2008}. Finally, measured transit times for each method were plotted against measured centerline distances along the aortic centerline. To estimate the relationship between distance and the time shift, simple linear regression was used to fit a line across the 3 data points, where the inverse of the fitted slope is the estimated PWV across all data points~\cite{wentlandReviewMRIbasedMeasurements2014}.

\begin{figure}[ht]
    \vspace{0.25cm}
    \centering
    \includegraphics[width=\textwidth]{time_shift.png}

    \caption[Time-shift methods]{Flow waveforms in the aorta were generated at multiple ROIs from time-resolved phase contrast images. Four different methods were used t
    o calculate PWV based on the time shift between flow curves.}\label{fig:time_shift}
    \vspace{0.25cm}
\end{figure}

\subsubsection{Analysis}
Aortic distance measurements, along with pulse wave velocities for each acquisition and time shift method, were tabulated for the entire cohort, by age category (50\-59 years, 60\-69 years, and 70\-79 years), and by sex. Mean, standard deviation, and number of data points (N) were reported for each measure. Note that non-physiological PWV values (e.g., negative PWV measures and PWV values above 20 m/s) were removed from analysis. For each time shift method, pulse wave velocities were compared between acquisitions using the Skillings-Mack test, a nonparametric repeated-measures ANOVA adjusted for missing values (analogous to the Friedman test). Additionally, box plots were used to visualize PWV distributions, outliers (defined using the interquartile rule), and group medians with line segments demonstrating paired measurements. Correlations between acquisition and time shift method pairs were computed using Pearson correlation coefficients (r) and were visualized using a correlation matrix. 

Simple linear regression was used to assess the effect of age on aortic pulse wave velocity for each acquisition and time shift method. For each model, regression plots were created and regression coefficients ($\beta$), coefficients of determination ($R^2$) and p-values were recorded. Statistical significance was defined as p < 0.05.


\section{Results}
Thirty-eight cases were excluded due to: poor plane prescription (N=8), insufficient image quality or velocity aliasing (N=11), or poor cardiac or respiratory gating (N=19) in one or more of the acquisition types, arriving at a final total of 45 cases. Of the 19 subjects that were excluded  Demographics for the analyzed study population are shown in Table 1.Thirty-eight cases were excluded due to: poor plane prescription (N=8), insufficient image quality or velocity aliasing (N=11), or poor cardiac or respiratory gating (N=19) in one or more of the acquisition types, arriving at a final total of 45 cases. Of the 19 subjects that were excluded  Demographics for the analyzed study population are shown in Table 1.

\begin{table}[H]
    \centering
    \caption[Patient Demographics]{Patient Demographics}
    \begin{tabular}{ll}
    \toprule
    Variable & Value \\
    \midrule
    Age (years) & $63.1 \pm 7.41$ (50--77) \\

    Sex & \\
    \quad Female & 34 (76\%) \\
    \quad Male & 11 (24\%) \\

    Race* & \\
    \quad White & 44 (97.8\%) \\
    \quad Unknown & 1 (2.2\%) \\

    Education (years) & $16.8 \pm 1.98$ (12--20) \\

    Parental history of dementia & \\
    \quad Yes & 36 (80.0\%) \\
    \quad No & 9 (20.0\%) \\

    Parental history of AD & \\
    \quad Yes & 23 (51.1\%) \\
    \quad No & 9 (20.0\%) \\
    \quad Unknown & 13 (28.9\%) \\

    APOE4 carrier** & \\
    \quad Yes & 30 (66.7\%) \\
    \quad No & 14 (31.1\%) \\
    \quad Unknown & 1 (2.2\%) \\
    \bottomrule
    \end{tabular}
    \begin{flushleft}
        \footnotesize
        For continuous data, values are expressed as mean ± SD and data ranges are in parenthesis. For categorical data, values are expressed as numerators with percentage in parenthesis. 
        *Race was self-reported and obtained as required by federal funding agencies.
        **APOE4 = apolipoprotein E4

    \end{flushleft}
\end{table}

\begin{figure}[ht]
    \vspace{0.25cm}
    \centering
    \includegraphics[width=\textwidth]{rad_images.png}

    \caption[Radial image examples]{Examples of time-averaged and time-resolved phase contrast images in the ascending aorta (AAo) and abdominal aorta (AbdAo) across a Cartesian product scan, a standard radial reconstruction, and a high-res (Locally Low Rank) radial reconstruction.}\label{fig:rad_images}
    \vspace{0.25cm}
\end{figure}

Average ± standard deviation temporal resolution was 26.2 ± 3.6 ms (dependent on heart rate) for the Cartesian reconstruction, 25.2 ± 3.6 ms for the standard radial reconstruction, and 12.5 ± 1.8 ms for the high temporal resolution (LLR) radial reconstruction. Plane distances (in the superoinferior direction) between the aortic arch and abdominal aorta 2D slices were 131 ± 18.3 mm. Representative images from each of the 3 acquisition methods are shown in Figure 4. Table 2 provides aortic lengths between each of the 3 measurement locations and PWV values for each acquisition and time shift method by the whole group, age category, and by sex.

\begin{table}[H]
    \centering
    \caption[Aortic Distances and Pulse Wave Velocities]{Aortic Distances and Pulse Wave Velocities}
    \small
    \begin{tabularx}{\textwidth}{llXXXXXX}
    \toprule
    Acquisition & Parameter & Overall & 50--59 & 60--69 & 70--79 & Female & Male \\
    \midrule

    Anatomical & Aortic Length & 
    $258.0 \pm 23.6$ (45) & 
    $252.1 \pm 21.1$ (13) & 
    $257.5 \pm 24.6$ (17) & 
    $263.6 \pm 24.7$ (15) & 
    $252.1 \pm 7.5$ (3) & 
    $264.3 \pm 23.1$ (13) \\

    \midrule
    \multicolumn{8}{l}{\textbf{Cartesian}} \\

    & TTP  & $7.2 \pm 2.3$ (45) & $6.5 \pm 1.2$ (13) & $6.6 \pm 2.0$ (17) & $8.4 \pm 2.7$ (15) & $7.5 \pm 1.2$ (3) & $6.9 \pm 2.2$ (13) \\
    & TTF  & $7.1 \pm 2.9$ (45) & $6.0 \pm 0.9$ (13) & $6.3 \pm 2.4$ (17) & $9.1 \pm 3.6$ (15) & $7.3 \pm 3.1$ (3) & $6.4 \pm 1.5$ (13) \\
    & TTU  & $7.3 \pm 2.8$ (45) & $6.2 \pm 1.0$ (13) & $6.5 \pm 2.4$ (17) & $9.1 \pm 3.4$ (15) & $7.5 \pm 2.2$ (3) & $6.6 \pm 1.7$ (13) \\
    & XCOR & $7.3 \pm 1.8$ (45) & $7.4 \pm 1.6$ (13) & $6.8 \pm 2.2$ (17) & $7.8 \pm 1.7$ (15) & $7.9 \pm 1.7$ (3) & $7.2 \pm 2.3$ (13) \\

    \midrule
    \multicolumn{8}{l}{\textbf{Radial \- Low Temporal Resolution}} \\

    & TTP  & $7.1 \pm 1.4$ (44) & $6.5 \pm 1.1$ (13) & $6.8 \pm 1.3$ (16) & $8.0 \pm 1.3$ (15) & $9.0 \pm 1.4$ (2) & $7.0 \pm 1.3$ (13) \\
    & TTF  & $7.0 \pm 2.0$ (44) & $6.2 \pm 1.2$ (13) & $6.2 \pm 1.4$ (16) & $8.6 \pm 2.1$ (15) & $10.5 \pm 3.5$ (2) & $6.8 \pm 1.0$ (13) \\
    & TTU  & $7.0 \pm 1.8$ (44) & $6.2 \pm 2.1$ (13) & $6.6 \pm 1.5$ (16) & $8.1 \pm 1.4$ (15) & $9.0 \pm 1.4$ (2) & $6.9 \pm 2.2$ (13) \\
    & XCOR & $7.2 \pm 2.1$ (44) & $7.2 \pm 2.0$ (13) & $6.9 \pm 2.1$ (16) & $7.4 \pm 2.2$ (15) & $10.5 \pm 2.1$ (2) & $7.2 \pm 2.3$ (13) \\

    \midrule
    \multicolumn{8}{l}{\textbf{Radial \- High Temporal Resolution}} \\

    & TTP  & $7.5 \pm 1.3$ (45) & $7.1 \pm 1.2$ (13) & $7.1 \pm 1.2$ (17) & $8.2 \pm 1.0$ (15) & $9.2 \pm 0.5$ (3) & $7.3 \pm 1.1$ (13) \\
    & TTF  & $7.3 \pm 1.6$ (45) & $6.6 \pm 1.2$ (13) & $6.9 \pm 1.5$ (17) & $8.4 \pm 1.5$ (15) & $9.4 \pm 1.1$ (3) & $7.0 \pm 0.9$ (13) \\
    & TTU  & $7.3 \pm 1.7$ (45) & $6.4 \pm 2.0$ (13) & $7.1 \pm 1.4$ (17) & $8.3 \pm 1.1$ (15) & $9.0 \pm 0.1$ (3) & $7.3 \pm 1.2$ (13) \\
    & XCOR & $7.6 \pm 2.0$ (45) & $7.5 \pm 2.1$ (13) & $7.4 \pm 2.1$ (17) & $7.8 \pm 1.9$ (15) & $9.7 \pm 3.3$ (3) & $7.8 \pm 1.9$ (13) \\

    \bottomrule
    \end{tabularx}
    \begin{flushleft}
        \footnotesize
        All values expressed as: mean ± SD (number of samples) \\

        PWV = pulse wave velocity; TTP = time-to-point; TTF = time-to-foot; TTU = time-to-upstroke; XCOR = cross correlation.


    \end{flushleft}
\end{table}


\begin{figure}[ht]
    \vspace{0.25cm}
    \centering
    \includegraphics[width=\textwidth]{boxplots.png}

    \caption[PWV distributions across methods]{Distributions of pulse wave velocities calculated using four different transit time methods across a Cartesian product scan, a standard radial reconstruction, and a high-res (Locally Low Rank) radial reconstruction. Connecting lines between paired data points highlight the individual correlations on a participant level.}\label{fig:boxplots}
    \vspace{0.25cm}
\end{figure}

Using the Skillings-Mack test, a significant difference in PWV across acquisition methods was observed with the TTP algorithm (Chi-Square = 8.62; P = 0.01) and the XCOR algorithm (Chi-Square = 7.20; P = 0.03); however no significant difference was observed for the TTF and TTU algorithms. Figure 5 shows box plots of PWV values for each MRI acquisition method. Qualitatively, it is apparent that Cartesian PWV values often do not align well with radial PWV values, particularly at higher PWVs. Pearson correlation coefficients between each time shift method-acquisition pair are shown as a correlation matrix in Figure 6A. All correlations between pairs were statistically significant. Mean correlations averaged across acquisitions (i.e., mean correlation in each transit time “block”) are shown in Figure 6B. The highest mean correlations were between TTF (mean r = 0.73) and XCOR (mean r = 0.72) measures compared across each acquisition. Comparing between transit time methods decreased mean correlation. Mean correlations averaged across transit time methods are shown in Figure 6C. The highest mean correlation between transit time methods was observed in the Cartesian acquisition (mean r =0.80); however, transit times were also highly correlated for the standard radial (mean r = 0.73) and high temporal resolution radial (mean r = 0.72) acquisitions. While comparing transit times between the two radial acquisitions showed high degree of correlation (mean r = 0.67), comparing transit times between the Cartesian and radial acquisitions showed poor correlation (mean r < 0.28). 

\begin{figure}[ht]
    \vspace{0.25cm}
    \centering
    \includegraphics[width=\textwidth]{correlation.png}

    \caption[PWV correlations]{A. Correlation matrix of Pearson correlation coefficients between individual pairs of transit time methods and acquisitions. The transit time methods include Time-To-Point (TTP), Time-To-Foot (TTF), Time-To-Upstroke (TTU), and Cross Correlation (XCOR) while the acquisitions include a Cartesian product sequence, a standard radial reconstruction, and a high temporal resolution (Locally Low Rank) radial reconstruction. B. Mean correlations averaged across acquisitions C. Mean correlations averaged across transit time methods}\label{fig:correlation}
    \vspace{0.25cm}
\end{figure}

For the standard and hi-res radial acquisitions, simple linear regression demonstrated a significant relationship between aortic PWV and age for the TTP, TTF, and TTU methods but no significant relationship using the XCOR method. The Cartesian acquisition demonstrated a significant relationship between aortic PWV and age for the TTF and TTU methods but not for TTP and XCOR\@. Regression plots for the TTF method are shown for all 3 acquisitions in Figure 7. Using the TTF method, regression coefficients ($\beta_{age}$) ranged from 0.06 to 0.11 m/s, R2 values ranged from 0.16 to 0.22, and standard errors ranged from 0.03 to 0.04. The high temporal resolution radial data showed the strongest relationship with age, with $R^2$ values ranging from 0.02 to 0.20 for all methods. 

\begin{figure}[ht]
    \vspace{0.25cm}
    \centering
    \includegraphics[width=\textwidth]{linear_regression.png}

    \caption[PWV linear regression with age]{Linear regression models showcasing the positive relationship between age and pulse wave velocity calculated via the Time-to-Foot method in three different acquisitions: a Cartesian product sequence, a standard radial reconstruction, and a high-res (Locally Low Rank) radial reconstruction.}\label{fig:linear_regression}
    \vspace{0.25cm}
\end{figure}

\subsection{Discussion}
In this work, we demonstrated the feasibility of a free-breathing, radial 2DPC (two-dimensional phase-contrast) MRI sequence for estimation of aortic pulse wave velocity (PWV). The free-breathing, radial acquisition was compared to the standard clinical breath-hold Cartesian sequence. From the radial acquisition, two datasets were reconstructed, one dataset with matched resolutions to the Cartesian acquisition and another with high temporal resolution obtained using an advanced local low rank reconstruction technique, leading to a total of 3 acquired MRI datasets. Four transit time methods were explored to investigate their impact on PWV\@. Mean PWV values were comparable between the 3 acquisitions and 4 time shift algorithms; average PWV across the entire cohort ranged from 7.1\-7.7 m/s. Nonparametric tests indicated that there were significant differences between acquisitions using the TTP and XCORR algorithms though not with the TTF and TTU algorithms. Lastly, PWV calculated with the TTF and TTU algorithms were strongly correlated with age for each acquisition method with the radial acqusitions additionally showing a correlation with the TTP method.

While PWV measures were comparable between the Cartesian, standard radial, and high temporal resolution radial datasets, there was a significant degree of pairwise variability. This is evident from the correlation matrices (Figure 6) and in the paired line segments in the box plots (Figure 5). It is still unclear what is causing this inconsistency. There are several potential contributors, including potential gating timing issues resulting from separate acquisitions, differences in heart rate between aortic arch and abdominal aorta scans, intra-acquisition heart rates (which may blur the waveforms), differences in flow waveform noise, respiratory effects, or underlying post-processing in the automatically-reconstructed Cartesian images (e.g., phase unwrapping, temporal interpolation, etc.). All acquisitions and transit time methods (aside from cross correlation) showed significant relationships with age, as expected. Specifically, we found that aortic PWV increases by around 1.1\-1.4 m/s each decade of life. This agrees closely with a large study (N=397) performed by van Hout et al., in which they assessed aortic PWV in older adults (mean age 55 ± 6 [SD] years) using 2DPC MRI~\cite{vanhoutNormalReferenceValues2021}. Using linear regression, they observed an average PWV increase of 0.9 m/s over 10 years. Similar results have been demonstrated in other 2DPC MRI studies~\cite{dyverfeldtQuantificationIntravoxelVelocity2006,parikhMeasurementPulseWave2016}. Furthermore, the study by van Hout reported a mean PWV value of 6.8 m/s using the TTF method for the age range 60\-65 years, which is similar to the results obtained here (mean PWV = 6.5 ± 2.4 m/s). Additionally, a large normative study (N=777, age ranges 21\-85 years) by Nethononda et al.\ reported a mean aortic PWV of around 9 m/s in the 60\-69 year age range~\cite{nethonondaGenderSpecificPatterns2015}. 

In regards to transit time algorithms, TTF was associated with low standard deviation, high correlation with other transit time algorithms, a significant correlation with age, and robustness (i.e., not prone to algorithm failure). While the time-to-point relies on a single point in the cardiac cycle and is thus susceptible to sampling errors and velocity noise, it was observed that the systolic upslope portion of the flow waveforms was often preserved even in noisier datasets. Cross correlation is beneficial in that the entire flow waveform is utilized to estimate time shifts, which theoretically reduces sampling errors. However, it was noted that XCOR algorithm tended to produce erroneous results if the systolic peak of the waveform was near the edge of the cardiac cycle. While TTF and XCOR showed promise in this study, many other studies utilized the TTP and TTU methods, including applanation tonometry. Thus, it may be beneficial to acquire all time shifts (which take relatively little time to compute) in order to allow for comparisons with other studies and to provide backup measures in the event one of the algorithms fails. 

The local low rank reconstruction employed in this study was used to increase the temporal resolution by a factor of 2. Despite severe undersampling (approximately 62 projections per frame with respiratory gating), image quality was preserved, and in some cases, improved relative to the standard 40 frame radial reconstruction. It was hypothesized that increased temporal resolution would lead to more accurate PWV estimations but the results show that PWV values obtained from these two reconstructions are similar. Interestingly, the high temporal resolution data showed the strongest correlation with age. Streaking artifacts were present in some datasets which appeared as noise in the obtained flow waveforms. While smoothing the waveforms reduced this artifact, velocity noise likely has an effect on the waveform shape and thus calculated pulse wave velocities. 

It should be noted that other phase contrast methods have been successfully employed to probe aortic stiffness, including in-plane 2DPC~\cite{westenbergImprovedAorticPulse2010}, “pencil beam” NMR~\cite{hardyPencilExcitationInterleaved1996}, the area-flow “QA” method~\cite{vulliemozEstimationLocalAortic2002}, real-time PWV~\cite{robertsRealtimeAorticPulse2015}, and more recently 4D flow MRI~\cite{dyverfeldtPulseWaveVelocity2014}. 4D flow MRI is a particularly promising application for several reasons: (1) it provides full coverage of the aorta in a single scan and allows for continuous flow measurements along the entire aorta, (2) it does not require tedious plane prescription for each slice, which is particularly useful in subjects with complex aortic shapes, (3) morphological data is also obtained during the scan and can be used to create an aortic centerline from the angiogram, (4) all flow and morphology data are spatially and temporally registered, (5) regional PWV measures can be calculated, and (6) other useful hemodynamic parameters can be computed. A major drawback of this method is the relatively poor temporal resolution, a result of multiple velocity encodings required to calculate velocity in each direction. However, recent advancements in acquisition and reconstruction have allowed for major improvements in spatial and temporal resolution~\cite{valvanoAccelerating4DFlow2017,rivera-riveraAssessmentVascularStiffness2021}. A recent study performed by Fico et al.\ demonstrated a correlation between aortic stiffness assessed using applanation tonometry and intracranial pulsatility derived from 4D flow MRI~\cite{ficoImpactAgingAssociation2022}. 

\subsubsection{Limitations}
There were several significant limitations to this study. While temporal resolutions were similar between the aortic arch and abdominal aorta scans, there were slight changes in heart between scans, resulting in abdominal aorta waveforms with different length (RR interval) and temporal resolution to varying degrees since the scans were not interleaved and were obtained separately. Additionally, because PPG gating was used, the systolic portion of the flow waveform was located in the last frames of the cine reconstructed 2DPC images. If there was a greater difference in nominal RR intervals (and thus temporal resolution), then accumulated time differences at the end of the waveform led to “shifting” of the part of the flow waveform used for transit time calculations. In order account for this effect, one set of time-stamps (from the aortic arch scan) were used for all waveforms such that the length of all waveforms were matched. This effect should be further studied using acquisition methods that obtain flow measurements simultaneously (e.g., simultaneous multi-slice, interleaved acquisitions, 4D flow MRI, etc.). It is also important to note that ECG gating may have different intrinsic and physiological time delays compared to PPG gating. This may be an important consideration for the prospectively-gated Cartesian acquisition, in which trigger window or trigger delays may effect waveform shape as discussed earlier. 

Additionally, due to retrospective gating, the radial acquisitions required longer scan times than for the Cartesian acquisitions. This may have allowed for greater changes in heart rate during the scan, which may lead to waveform blurring. Differences in respiratory phase may also cause slices to shift vertically even though an identical plane prescription location was used for both radial and Cartesian scans. Furthermore, even though radial data is respiratory-gated during a specific phase, motion may still exist within the acceptance window which may introduce slight image blurring. The physiological effects of respiratory phase (e.g., forced breath-hold compared to free-breathing state) may also introduce subtle biases between acquisitions. 

\subsubsection{Future Directions}
Elongation and increased tortuosity of the aorta has been observed to change with age and has been hypothesized to be a result of weakened vessel wall structure~\cite{sugawaraAgeassociatedElongationAscending2008}. Future studies could quantitatively measure these morphological characteristics and assess their correlations with age and PWV measures. Additionally, differences in stiffness along the aorta have been reported in the literature; specifically, the proximal aorta tends to be the most elastic and stiffens distally. Because time shifts and distance measures are calculated at three points along the aorta, measures of proximal aorta and abdominal aorta stiffness could be separately calculated at the expense of measurement noise. This may be of particular relevance to a relatively new biomarker, termed PWV ratio, which calculates a stiffness gradient along the arterial tree and may help elucidate the degree of impedance mismatching that can result in higher pulse pressure transmission to the microcirculation~\cite{covicPulseWaveVelocity2015}. Other areas of investigation may include evaluating the effect of pre- and post-menopausal changes within our cohort. Studies have suggested a potential hormonal effect that may differentially affect vessel wall stiffness in women after menopause which could be assessed with the present cohort, consisting of 34 females~\cite{nethonondaGenderSpecificPatterns2015}. A recent study using 4D flow MRI cranial scans obtained internal carotid PWV measures using local low rank reconstructions in the WRAP and ADRC cohorts~\cite{rivera-riveraAssessmentVascularStiffness2021}. This same methodology could be applied here, as cranial 4D flow MRI is included in the LIFE MRI protocol. An interesting future direction would be assess the relationship between aortic and internal carotid artery PWV\@. 

Future technical improvements could include the application of machine learning in segmenting and skeletonizing the aorta for automatically generated aortic centerlines for distance measurements. Automatic segmentation of the aorta in the 2DPC images should also be explored to decrease processing time, increase accuracy of vessel area measurements, and increase repeatability. This would also allow for an additional measure of local PWV (“QA” method), that requires accurate flow and area quantification34. 

\subsection{Conclusion}
In conclusion, the radial, free-breathing 2DPC MRI sequence offers an alternative imaging technique for assessing aortic stiffness which may be particularly beneficial for subjects with breath-hold difficulty. The radial sequence offers a great deal of flexibility for physiological gating and for image reconstruction, which allowed us to generate a high temporal resolution dataset using local low rank reconstruction. Mean pulse wave velocities were comparable between the radial and Cartesian acquisitions and agree are externally consistent with other studies, however, the correlation between paired measures is relatively poor. Pulse wave velocities were tabulated in our cohort of 83 cognitively unimpaired older adults and showed strong associations with age. Future studies should seek to study intracranial hemodynamics alongside the obtained aortic pulse wave velocity data.
